\documentclass[12pt,letterpaper]{amsart}
\usepackage[T1]{fontenc}
\usepackage{newtxtext,newtxmath}
\usepackage[letterpaper,margin=1in]{geometry}
\usepackage{microtype}
\usepackage{booktabs,longtable,array}
\usepackage[hidelinks]{hyperref}
\title{Revision log for \emph{On twisted orbital-integral identities for $PGL(3)$ over a $p$-adic field}}
\author{David Joyner and ChatGPT (OpenAI)}
\date{Corrected revisit, Draft 1 --- November 23, 2026}
\begin{document}
\begin{abstract}
This AI-aided editorial and mathematical audit accompanies Draft 1 of the
corrected revisit, with mathematical contributions from ChatGPT (OpenAI).
It distinguishes inherited claims, newly proved elementary statements,
conditional reductions, and unverified historical proof obligations.
\end{abstract}
\maketitle

\section*{Typesetting corrections (November 23, 2026)}
The table of contents now uses separate indentation and number-label
widths for sections and subsections, including the unnumbered end
sections. The subtitle is typeset in a smaller italic font, separately
from the uppercase main title. Underlining has been removed from
bibliography entries: book titles are italicized and journal volume
numbers are bold. The title-page note now describes the paper more concisely as a significant geometric rewrite of~\cite{Jo90}. In addition, several new TikZ diagrams were added to the main manuscript to illustrate apartment folding, the Levi subcomplex $\mathcal Z$, midpoint fixed points, and tree parity. No mathematical statements or proofs were changed by these presentational edits.

\section*{Structural reorganization (November 23, 2026)}
The stand-alone notation section has been moved earlier and split into two subsections of Section~1: ``Notation'' and ``Conventions.'' In addition, the historical sections formerly appearing later as ``The buildings for $PGL(3)$ and $SL(2)$'' and ``The buildings for $PGL(3)$ and $PGL(2)$'' have been relocated as subsections of the general background section. Their labels and historical equation tags were preserved.

\section*{Additional geometric figures (November 23, 2026)}
Several further TikZ diagrams were added: a local $\widetilde A_2$ star at a hyperspecial vertex, a flag--chamber correspondence figure for the factor $Q$, and a schematic comparison of vertex-centered and edge-midpoint-centered balls in the $PGL_2$ tree.

\section*{Expanded introduction and status of shorter sections (November 23, 2026)}
The new introduction retains the early Subsections~1.1 and~1.2 for
Notation and Conventions. It now adds distinct subsections explaining
the historical strategy and provenance, orthogonal duality and the
Levi subcomplex, lattice types and ramification, the finite-flag
normalization $Q$, tree-Hecke recurrences and Satake theory, the
conditional signed tree count, and the status of the shorter sections.
The old single-paragraph section ``Outline of the original paper''
has been incorporated into the historical-strategy subsection of
the introduction. The previously separate historical section named
``Introduction'' is retitled ``Historical framework: norm maps and fundamental lemmas''
to distinguish it from the new, expanded introduction. The wording was revised to distinguish the
attributions of the 1990 work from new claims of the revisit. The
historical building and orbital-integral sections and all their
existing statement labels were retained. Subsequent numbered
sections shift down by one, but internal references use labels.

\emph{Mathematical verification status is unchanged.}
The short parity, flag, and graph-operator sections are compact
proofs, not placeholders. The characteristic-leaf identifications,
complete transfer-factor sign, measure audit, and any independent
Hecke proof of the split identities remain unproved here. The
introduction points explicitly to the essential To Do items rather
than treating these gaps as completed research. This change is
expository and structural; it introduces no new theorems.


\section{Materials read and provenance}
The uploaded corrected Draft 2 LaTeX source, its accompanying Markdown
revision log, and the author's mathematical paper style guide were
read in full. Their filenames and original versions are preserved
with the user's submitted materials.
The active GitHub connector confirmed that
\href{https://github.com/wdjoyner/trace-formulas}{the trace-formulas repository} is a repository
accessible to the author, but its \texttt{main} branch did not have an
\path{adjoint-lift} directory during this session. The requested
\path{adjoint_lift_combined.tex} and
\path{adjoint_lift_referee_report.tex} therefore could not be
read. They were not reconstructed from memory or invented.
The 1990 publication was located on the \emph{Canadian Journal of
Mathematics} site and its title-page publication data and selected
mathematical passages were inspected. This was not an exhaustive
page-by-page collation.

The original historical source files do not correspond exactly to the
published article, as specified by David Joyner. Therefore this
working draft is neither a facsimile nor a literal transcription.
David Joyner states that the original work was done during official
federal duties and is public-domain United States government work.
The new 2026 manuscript is separately identified as a corrected revisit.

\section{Corrected Draft 2 results retained}
The following corrections were already made in the uploaded Draft 2,
and were carried forward without claiming a second independent
verification of every geometric lemma:
\begin{enumerate}
\item The representation is $\operatorname{Sym}^{2n}(\mathrm{Std})\otimes
(\det)^{-n}$, and the $n=1$ adjoint matrices and basis were corrected.
\item The normalized building representative is $A=CJ$ for the actual
outer involution $\sigma_J(C)=JC^{-t}J$.
\item Twisted stabilizers, quotient integration variables, and rational
stable-class kernels were corrected. The quadratic-torus case admits
the stated cohomological simplification.
\item The $\lambda_2$ Satake character restriction is $k_1-k_2$.
The full Weyl discriminant and the nonabelian Jacobian interpretation
were restored.
\item $\mathcal C(\Lambda)$ (a double coset) and $Q(\Lambda)$ (a
scalar) were separated. An open finite-index subgroup $U\subset F^\times$
replaced the incompatible compact finite-index description.
\item The contradictory parity wording of the old Lemma 4.14 and the
fixed-vertex notation of the old formula (4.19) were corrected.
\item The Section 5 scalar relation was restored as
$S(0)/Q=1$, equivalent to $S(0)=Q=s(0)$ under its hypotheses.
\item The ramified case $|1+\beta|_E<1$ was retained and illustrated
by $F=\mathbb Q_3$, $E=F(\sqrt 3)$, $\beta=2+\sqrt3$.
\item The signed comparison was explicitly distinguished from the
absolute-value result $|S(0)|=s(0)$.
\end{enumerate}
These items reflect the uploaded revision log. They are not described
as new discoveries of the present revisit.

\section{New results established in Draft 1}
The results below were proved from definitions by explicit elementary
arguments for this draft (verification Level 1). Their novelty is
\emph{expository or organizational}: they are not claimed to be
previously unknown theorems of building theory or harmonic analysis.
\begin{longtable}{@{}p{0.38\linewidth}p{0.53\linewidth}@{}}
\toprule
Result in the new paper & Verification and purpose\\
\midrule
\endhead
Apartment folding & The action of $\sigma_J$ on diagonal apartment
coordinates fixes the line $x_1+x_3=2x_2$. Fixed integral vertices
have exponent vectors $(m,0,-m)$. Direct calculation.\\
Diagonal midpoint obstruction & The involution
$[L]\mapsto[A L^\vee]$, $A=\operatorname{diag}(\pi^{d_i})$,
fixes the apartment point $\boldsymbol d/2$. A vertex exists
in this apartment iff all $d_i$ agree modulo two.\\
Type congruence & For $PGL_n$, a twisted fixed vertex satisfies
$2t_n([L])\equiv v_F(\det A)\pmod n$. Direct determinant
valuation; necessary, not sufficient.\\
Tree parity & Odd determinant valuation exchanges the two
vertex types on the $PGL_2$ tree. A fixed geometric point,
if present, is an inverted edge midpoint. Explicit examples.\\
Flag/Iwahori factor & $[K:I]=(q^2+q+1)(q+1)$ and
$Q=q^{-3}[K:I]=\sum_{w\in S_3}q^{-\ell(w)}$.
Exact finite projective incidence and Weyl-length proof.\\
Iwahori averaging & With $\operatorname{vol}(K)=1$,
$e_K=[K:I]^{-1}\sum_{kI\in K/I}{}^ke_I$ as functions.
Direct finite-coset proof.\\
Rank-two paths & $U_1^2=(q+1)U_2+U_{(2,0,0)}$.
Elementary-divisor classification of two-step lattice chains.\\
Tree distance recurrence & $AD_1=D_2+(q+1)D_0$ and
$AD_n=D_{n+1}+qD_{n-1}$ for $n\geq2$; operator generating
series $(1-u^2)/(1-Au+qu^2)$. Exact graph proof.\\
Adjoint Satake character & $\operatorname{tr}(\operatorname{Sym}^2
\operatorname{diag}(z,z^{-1}))=z^2+1+z^{-2}$. Direct weight proof.\\
Tree-ball difference & $B_v(m)-B_e(m)=q^m$ with explicit
rational generating functions. Complete counting proof and
numerical examples at $q=3,5$.\\
Pairing criterion & If the extra signed contributions admit a
weight-preserving bijection, their difference is the base-ball
difference $q^m$. Complete proof of the \emph{conditional} statement.\\
\bottomrule
\end{longtable}

\section{Cited results distinguished from new proofs}
The assertion that the geometric fixed set of the orthogonal
involution equals the building of its connected fixed group is
attributed to Prasad--Yu, \emph{Invent. Math.} \textbf{147}
(2002), Theorem 1.9. The central hypotheses are that the finite
involution group has order prime to the residue characteristic and
that the ambient group is reductive. For $PGL_3$ and its split
orthogonal involution, odd residue characteristic satisfies the
prime-to-$p$ condition. This is a cited result (Level 2 for
publication and theorem identity, based on the accessible original
introduction and indexed theorem references); it is not an
independent proof of Prasad--Yu. The author should consult the
complete original theorem text before submission.

The identifications involving quadratic norm characters and the
Satake algebra are standard facts of local fields and spherical
harmonic analysis. Their elementary special cases are derived
where feasible. Any larger endoscopic assertion is cited or
retained as a conditional historical input.

\section{Unresolved issues and precise failure points}
\begin{enumerate}
\item \textbf{Transfer factor sign: essential.} The sign in the
historical formula (3.18) has not been determined by this revisit.
The Hilbert norm character does not fix an independent initial sign.
The old Proposition 5.6 yields an absolute-value comparison only.
This is a genuine missing verification, not a proved sign choice.
\item \textbf{Building geometry: essential.} The historical Lemmas
4.9, 4.13, 4.16, 4.18, 5.3, and 5.4 are retained, but the fixed-vertex
and characteristic-leaf counts have not been independently reconstructed.
The difficult boundary cases $b_1-b_2=-1,0$ in Lemma 5.4 remain.
\item \textbf{Ramified vanishing: essential.} The last vanishing
argument uses the preceding lemmas and passage to subdivided trees.
It was not independently certified. The explicit $\mathbb Q_3$
example shows that ramified parameters satisfying the supposedly
vanishing condition really exist.
\item \textbf{Common correction $R$: essential.} The proved identity
$B_v(m)-B_e(m)=q^m$ only computes the base counts. A
weight-preserving pairing of extra characteristic leaves has not
been constructed, so equality of the correction terms is a hypothesis,
not a theorem of the revisit.
\item \textbf{Haar and nonunit conventions: essential for generalization.}
The $q^{k_3-k_1}$ factor and $Q(\Lambda)$ are inherited from the
historical source. They require a complete audit of double-coset
representatives, measures, and Macdonald's formula outside $\Lambda=0$.
\item \textbf{Prasad--Yu in residual characteristic two: excluded.}
The present application assumes odd residue characteristic.
Nothing is claimed for wild or residue-characteristic-two cases.
\end{enumerate}

\section{Expository and structural changes}
The original five numbered mathematical sections and their theorem
labels were retained as the historical core. The old opening
was replaced by a new abstract, motivation, two-period literature
review, explanatory background, and notation table. Five substantial
sections add proofs and examples concerning folding, parity,
finite projective flags, Hecke adjacency, and signed tree counts.
There is a synthesis section, glossary, and an extended To Do section
immediately before acknowledgments and references.
Historical discussion of detected errors belongs in this log,
not in the new mathematical proofs. The inherited original equations
retain many explicitly tagged historical equation numbers; this is
an intentional preservation of their cross-reference concordance.
All new equations are numbered automatically.

\section{Concordance and renumbering}
The inherited historical labels \texttt{sec:1} through
\texttt{sec:5}, together with \texttt{lemma:4.5},
\texttt{lemma:4.16}, \texttt{lemma:5.4}, and
\texttt{prop:5.6}, were not renamed. Since new material follows
those sections, their internal historical numbering remains
usable. The unnumbered preliminaries and afterword are newly
organized. New results carry labels beginning
\texttt{prop:new-} or \texttt{eq:new-}; there is no
theorem-number concordance needed for those statements.

The historical reference keys were changed to author-year keys in
the new compilation. The exception \texttt{LaUn} denotes Langlands'
unpublished notes, whose exact bibliographic date has not been
independently established. The new citations
\texttt{Jo90}, \texttt{PY02}, and \texttt{Pr18} were added.

\section{AI contributions and verification limitations}
OpenAI ChatGPT, listed as a co-author at David Joyner's explicit
request for this working draft, wrote the explanatory sections,
constructed and checked the direct lattice, incidence, and graph
proofs, developed the conditional pairing formulation, and produced
this source and log. It located bibliographic entries and compared
selected details with web-accessible source descriptions. It did
\emph{not} independently prove Langlands' omitted geometric lemmas,
verify the historical article page by page, or compute the
transfer-factor sign. AI mathematical review is not independent
certification. For submission to a journal or arXiv, authorship
and AI-disclosure requirements must be checked separately.

\section{Compilation and distribution}
The four deliverables are the editable new manuscript, its compiled
PDF, the separate LaTeX revision log, and its compiled PDF.
The compilation is with pdfLaTeX, 12-point Times-style text
and mathematics via \texttt{newtx}, one-inch margins,
and US letter paper. The author may place these files under
\path{trace-formulas/adjoint-lift/} after reviewing the
proof status, with the new manuscript distinguished from the
original combined file and the corrected historical Draft 2.

\begin{thebibliography}{9}
\bibitem[Jo90]{Jo90}
D. Joyner, ``On twisted orbital integral identities for $PGL(3)$ over
a $p$-adic field,'' \emph{Canad. J. Math.} \textbf{42} (1990),
1098--1130, \url{https://doi.org/10.4153/CJM-1990-059-6}.
\bibitem[PY02]{PY02}
G. Prasad and J.-K. Yu, ``On finite group actions on reductive groups
and buildings,'' \emph{Invent. Math.} \textbf{147} (2002), 545--560.
\bibitem[KS99]{KS99}
R. Kottwitz and D. Shelstad, \emph{Foundations of Twisted Endoscopy},
Ast\'erisque \textbf{255} (1999).
\end{thebibliography}
\end{document}
